\documentclass[11pt]{article}

\usepackage[margin=1in]{geometry}
\usepackage{amsmath,amssymb}
\usepackage{enumitem}
\usepackage{hyperref}
\usepackage{natbib}

\newcommand{\langref}[1]{\textbf{#1}}
\newcommand{\langfam}[2]{\textbf{#1$_{#2}$}}

\title{Transformation Notes}
\date{\today}

\begin{document}
\maketitle

\section{Purpose}
This note records a reusable proof pattern for negative transformation results.
The key point is that many unsupported transformation claims do not need a bespoke lower-bound construction: they follow from succinctness lower bounds plus a trivial embedding of clauses or terms into the target language.

\section{General Meta-Rules}
\subsection{$\land$C from CNF lower bounds}
Suppose a language $L$ satisfies both of the following:
\begin{enumerate}[noitemsep]
  \item Every clause can be converted into an equivalent sentence of $L$ in polynomial time and polynomial size.
  \item \langref{CNF} does not compile to $L$ in polynomial time.
\end{enumerate}
Then $L$ cannot satisfy $\land$C in polynomial time.

The proof format used in the table is:
\begin{quote}
Any clause can be represented in $L$ because [construction]. If $L$ supported $\land$C, then CNF could compile into $L$ by conjoining all clauses. Since CNF cannot compile to $L$ [in polynomial / quasipolynomial time], this is impossible.
\end{quote}

\subsection{$\lor$C from DNF lower bounds}
Suppose a language $L$ satisfies both of the following:
\begin{enumerate}[noitemsep]
  \item Every term can be converted into an equivalent sentence of $L$ in polynomial time and polynomial size.
  \item \langref{DNF} does not compile to $L$ in polynomial time.
\end{enumerate}
Then $L$ cannot satisfy $\lor$C in polynomial time.

The proof format used in the table is:
\begin{quote}
Any term can be represented in $L$ because [construction]. If $L$ supported $\lor$C, then DNF could compile into $L$ by disjoining all terms. Since DNF cannot compile to $L$ [in polynomial / quasipolynomial time], this is impossible.
\end{quote}

\subsection{Quasipolynomial variants}
The same argument works with ``not in quasipolynomial time'' in place of ``not in polynomial time.''
It also works when the lower bound is obtained by propagation rather than from a primary source directly, as long as the propagated succinctness claim is already accepted as sound.

\section{Small Clause Constructions}
This section records explicit polynomial-size clause embeddings for languages that may need direct $\land$C lower bounds.

\subsection{Decision-diagram families}
Let $C = \ell_1 \lor \cdots \lor \ell_k$ be a clause.

\paragraph{\langref{FBDD}, \langref{nFBDD}, \langref{uFBDD}.}
Build a read-once decision chain on the variables of the literals in $C$.
At the node for literal $\ell_i$:
\begin{itemize}[noitemsep]
  \item if $\ell_i$ is positive, the $x_i=1$ edge goes to the 1-sink and the $x_i=0$ edge continues;
  \item if $\ell_i$ is negative, the $x_i=0$ edge goes to the 1-sink and the $x_i=1$ edge continues.
\end{itemize}
The final continuation edge goes to the 0-sink.
This is an FBDD of size $O(k)$, hence also an nFBDD and a uFBDD.

\paragraph{\langref{OBDD}, \langfam{OBDD}{<}, \langref{nOBDD}, \langfam{nOBDD}{<}, \langref{uOBDD}, \langfam{uOBDD}{<}.}
Use the same read-once decision chain, but follow the global order $<$.
Variables not appearing in the clause are simply skipped.
This yields an OBDD$_<$ of size $O(k)$.
Therefore the same clause is also small in OBDD, nOBDD, nOBDD$_<$, uOBDD, and uOBDD$_<$.

\subsection{DNNF-style families}
\paragraph{\langref{dec-DNNF}, \langref{dec-SDNNF}, \langfam{dec-SDNNF}{T}.}
Use the OBDD/OBDD$_<$ chain above and invoke the known polynomial compilations
\[
\langref{OBDD} \to \langref{dec\mbox{-}SDNNF},
\qquad
\langfam{OBDD}{<} \to \langfam{dec\mbox{-}SDNNF}{T}.
\]
For dec-DNNF, use the polynomial compilation \langref{FBDD} $\to$ \langref{dec-DNNF}.
Since the source clause representation is linear-size, the target representation is polynomial-size.

\paragraph{\langfam{d-SDNNF}{T}, \langfam{SDNNF}{T}.}
Fix a right-linear vtree $T$ induced by an order $<$.
Start from the OBDD$_<$ chain for the clause.
Compile it to \langfam{dec-SDNNF}{T}; this is automatically a \langfam{d-SDNNF}{T} sentence, and hence also an \langfam{SDNNF}{T} sentence.

\paragraph{\langref{SDD}, \langref{cSDD}.}
Start from an OBDD representation of the clause.
Every OBDD is a right-linear SDD, and in fact already compressed, so this gives polynomial-size representations in both SDD and cSDD. \citep{Darwiche_2011}

\section{Small Term Constructions}
This section records explicit polynomial-size term embeddings for languages that may need direct $\lor$C lower bounds.

\subsection{Decision-diagram families}
Let $\gamma = \ell_1 \land \cdots \land \ell_k$ be a term.

\paragraph{\langref{FBDD}, \langref{nFBDD}, \langref{uFBDD}.}
Build a read-once decision chain on the literals of $\gamma$.
At the node for literal $\ell_i$:
\begin{itemize}[noitemsep]
  \item if $\ell_i$ is positive, the $x_i=1$ edge continues and the $x_i=0$ edge goes to the 0-sink;
  \item if $\ell_i$ is negative, the $x_i=0$ edge continues and the $x_i=1$ edge goes to the 0-sink.
\end{itemize}
The final continuation edge goes to the 1-sink.
This is an FBDD of size $O(k)$, hence also an nFBDD and a uFBDD.

\paragraph{\langref{OBDD}, \langfam{OBDD}{<}, \langref{nOBDD}, \langfam{nOBDD}{<}, \langref{uOBDD}, \langfam{uOBDD}{<}.}
Use the same chain while respecting the global order $<$.
Variables absent from the term are skipped.
This gives an OBDD$_<$ of size $O(k)$, so the same term is small in every ordered, nondeterministic ordered, or unambiguous ordered variant.

\subsection{DNNF-style families}
\paragraph{\langref{dec-DNNF}, \langref{dec-SDNNF}, \langfam{dec-SDNNF}{T}.}
Use the same source constructions as above.
An OBDD term embeds into dec-SDNNF, and an OBDD$_<$ term embeds into dec-SDNNF$_T$.
For dec-DNNF, an FBDD term embeds via the polynomial compilation FBDD $\to$ dec-DNNF.

\paragraph{\langfam{d-SDNNF}{T}.}
Fix a right-linear vtree $T$ and build the OBDD$_<$ term chain.
Compile to dec-SDNNF$_T$ and regard the result as a d-SDNNF$_T$.

\paragraph{\langref{SDD}, \langref{cSDD}.}
Every OBDD term chain is already a right-linear SDD and a compressed right-linear SDD.
Hence terms are polynomial-size in both SDD and cSDD.

\section{How the Constructions Extend to Each Language}
The proof obligations split cleanly:
\begin{enumerate}[noitemsep]
  \item show a clause/term is polynomial-size in a base BDD-like language;
  \item use a known polynomial compilation into the target language;
  \item combine this with a propagated CNF or DNF lower bound for the target language.
\end{enumerate}

This route applies cleanly to the following families on the conjunction side:
\begin{itemize}[noitemsep]
  \item \langref{nFBDD}, \langref{uFBDD}: inherit small clauses/terms directly from FBDD;
  \item \langref{nOBDD}, \langfam{nOBDD}{<}, \langref{uOBDD}, \langfam{uOBDD}{<}: inherit small clauses/terms directly from OBDD$_<$;
  \item \langref{dec-DNNF}: obtain small clauses/terms by compiling FBDD chains;
  \item \langref{dec-SDNNF}, \langfam{dec-SDNNF}{T}: obtain small clauses/terms by compiling OBDD / OBDD$_<$ chains;
  \item \langfam{d-SDNNF}{T}, \langfam{SDNNF}{T}: use right-linear vtrees and the dec-SDNNF$_T$ embedding;
  \item \langref{SDD}, \langref{cSDD}: use the fact that OBDDs are right-linear SDDs, and already compressed in the canonical case;
  \item \langfam{SDD}{T}, \langfam{cSDD}{T}: on a right-linear vtree, these are exactly the OBDD / canonical OBDD fragments, so the same clause and term chains apply directly.
\end{itemize}

For DNNF and SDNNF, the same route also works:
\begin{itemize}[noitemsep]
  \item a clause is already a DNNF sentence, so DNNF gets a direct unconditional $\land$C lower bound from the propagated CNF $\not\leq$ DNNF gap;
  \item a clause has a small OBDD$_<$ representation, and OBDD$_<$ embeds into a right-linear SDNNF, so SDNNF gets the same kind of unconditional $\land$C lower bound from CNF $\not\leq$ SDNNF.
\end{itemize}

The same strengthened wording also applies to fixed-vtree structured variants when the quasipolynomial succinctness gap is already present:
\begin{itemize}[noitemsep]
  \item \langfam{SDD}{T} and \langfam{cSDD}{T} get direct unconditional $\land$C and $\lor$C lower bounds from the propagated CNF and DNF gaps into those languages;
  \item \langfam{d-SDNNF}{T} gets a direct unconditional $\land$C lower bound from the propagated CNF $\not\leq$ \langfam{d-SDNNF}{T} gap.
\end{itemize}

On the disjunction side, the route only applies when DNF does not already compile efficiently into the language.
In particular:
\begin{itemize}[noitemsep]
  \item it \emph{does} apply to dec-DNNF, dec-SDNNF, dec-SDNNF$_T$, d-SDNNF$_T$, uOBDD, uOBDD$_<$, SDD$_T$, cSDD$_T$, and cSDD;
  \item it \emph{does not} apply to nFBDD, SDNNF, nOBDD, or nOBDD$_<$, because DNF already compiles to those languages in polynomial time.
\end{itemize}

For the fixed-vtree families, the upgraded statements are genuinely quasipolynomial lower bounds, not merely polynomial lower bounds: the needed succinctness gaps are already present in the map as
\[
\langref{CNF} \not\leq \langfam{SDD}{T},\quad
\langref{CNF} \not\leq \langfam{cSDD}{T},\quad
\langref{CNF} \not\leq \langfam{d\mbox{-}SDNNF}{T},
\]
and similarly on the DNF side for $\lor$C.

\section{Pending Conditional Results}
Some results are not yet added because the best current succinctness premise is conditional or only excludes polynomial-time compilation:
\begin{itemize}[noitemsep]
  \item \langref{d-DNNF}: the natural $\lor$C argument should give a conditional ``not in polynomial time; quasipolynomial open'' result, but this is already present through propagation.
  \item \langref{uFBDD}: the same style of argument appears available, but is being deferred until the underlying DNF lower bound becomes unconditional.
\end{itemize}

\section{Current Scope}
This note is only about the compilation-based meta-rules for $\land$C and $\lor$C.
It does not attempt to systematize the more intricate selector-variable argument for FO.

\bibliographystyle{plainnat}
\bibliography{refs}

\end{document}
